\section{Data Collection}
\label{chap:data_collection}
In this chapter, we discuss the methods used to collect data during the project’s testing phase.

\subsection{Data Categories}

In chapter \ref{chap:rqm} we briefly discussed our research goal and questions: these questions study player engagement, motivation, enjoyment, and habits. Consequently, we chose several methods for collecting data to answer these questions. The data we gather all speaks to the game’s social benefit because a player’s engagement, motivation, enjoyment, and habits affect the individual and affect the player’s surroundings.

\textbf{Player engagement}\\
Data about the player’s engagement was gathered to see how players engage in the game during the two weeks. For this data, it was interesting to see how game metrics and other players’ presence affect the player’s engagement in the app. This data coincides with RQ1 and RQ4.

\textbf{RQ1:} \textit{How do game metrics affect the players’ engagement in our game concept?}\\
\textbf{RQ4:} \textit{How does social interaction affect the players’ engagement in our game concept?}

\textbf{Player enjoyment}\\
Data on player’s enjoyment was collected to investigate what makes the game enjoyable. Theory on what makes games enjoyable has been applied to the concept; hence, it was interesting to see how the game metrics affected the player’s enjoyment. Additionally, it was interesting to see how the game’s social aspects affect the player’s enjoyment. Data on player enjoyment will be used to reason about RQ2 and RQ5. 

\textbf{RQ2:} \textit{How do game metrics affect the players’ enjoyment in our game concept?} \\
\textbf{RQ5:} \textit{How does social interaction affect the players’ enjoyment in our game concept?}

\textbf{Player motivation}\\
To collect data on the player’s motivation, we choose to investigate how the game metrics and presence of other players affect the player’s motivation and willingness to be active and do activities outside their home. The data collected on this matter lay the foundation for reasoning about RQ3 and RQ6.

\textbf{RQ3:} \textit{How do game metrics affect the players’ motivation in our game concept?} \\
\textbf{RQ6:} \textit{How does social interaction affect the players’ motivation in our game concept?}

\textbf{Player habits}\\
As the game aims to be socially beneficial, it is interesting to investigate how it affects the player’s habits of being active outside their home. Data on this matter relates to RQ7.

\textbf{RQ7:} \textit{How does the use of our game concept affect the players’ habits?}

\subsection{Questionnaire}
In order to collect quantitative data, the participants were asked to answer a questionnaire. The questionnaire included a set of questions and statements that lay the foundation to reason about our research questions. The participants were given the questions in Norwegian (see complete questionnaire in Appendix \ref{appendix:questionnaire}). The questionnaire comprised 32 statements and 26 questions grouped into five sections depending on what type of data they could reason about.

We chose a 5-point Likert scale for the statements, where the respondents had the alternatives: strongly disagree, disagree, neutral, agree, and strongly agree

\textbf{General questions for the participant}\\
To gather data about participant groups, we initially collected some demographic data and data on their overall experience and use of the game. Additionally, we ended the questionnaire with a general feedback field for collecting relevant data that the questionnaire did not cover, such as feedback on how the testing was conducted, the app, and the concept. These questions are found in Table \ref{tab:general_questions}.


\begin{table}[H]
    \centering
    \begin{tabularx}{0.95\textwidth}{|c|X|}
        \hline
        \textbf{ID} & \multicolumn{1}{c|}{\textbf{Questions}} \\ \hline
        \textbf{Q1} & What is your gender? \\
        \textbf{Q2} & Are you a student \\
        \textbf{Q3} & What is your age? \\
        \textbf{Q4} & What is your email address? \\
        \textbf{Q5} & What is your username in Hover? \\
        \textbf{Q6} & Did you use the application at least once during the test period? \\
        \textbf{Q7} & Did you experience any technical challenges during the test period? \\
        \textbf{Q8} & Have you used Hold, Strava, or Pokémon Go before? \\ 
        \textbf{Q9} & Does your phone run iOS or Android? \\ 
        \textbf{} & ...  \\ 
        \textbf{Q26} & Do you have any other feedback?  \\ \hline
    \end{tabularx}
    \caption{General questions}
    \label{tab:general_questions}
\end{table}

\subsubsection{Player Engagement}
The subjects were given six statements and six questions to collect data on the player’s perceived engagement. The statement was based on the game metrics and social features that the game includes. The purpose of the statements was to reason about the perceived effect these metrics had on the player’s engagement. Table \ref{tab:player_engagement_statements} present the statements related to player engagement.

\begin{table}[H]
    \centering
    \begin{tabularx}{0.95\textwidth}{|c|X|}
        \hline
        \textbf{ID} & \multicolumn{1}{c|}{\textbf{Statements}} \\ \hline
        \textbf{S1} & I was curious about what medals I could achieve during the game. \\
        \textbf{S2} & I received new medals and rewards throughout the test period. \\
        \textbf{S3} & The game had sufficient challenge. \\
        \textbf{S4} & Challenges between friends inspired me to use the application more. \\
        \textbf{S5} & I often reacted to other players' activities/medals. \\
        \textbf{S6} & The notifications made me use the application more. \\ \hline
    \end{tabularx}
    \caption{Statements related to players' engagement.}
    \label{tab:player_engagement_statements}
\end{table}

Furthermore, we asked the subjects to quantify their use of the game and particular features. These questions can indicate how the metrics trigger players’ competitiveness and willingness to engage in the game. The questions on player engagement are presented in Table \ref{tab:player_engagement_questions}

\begin{table}[H]
    \centering
    \begin{tabularx}{0.95\textwidth}{|c|X|}
        \hline
        \textbf{ID} & \multicolumn{1}{c|}{\textbf{Questions}} \\ \hline
        \textbf{Q10} & How many of the days during the test period did you use Hover? \\
        \textbf{Q11} & On average, how long (number of minutes) did you use Hover every day during the test period? \\
        \textbf{Q12} & How many times during the test period did you accept challenges from friends? \\
        \textbf{Q13} & How many times during the test period did you challenge friends? \\
        \textbf{Q14} & On average, how many reactions did you give per day during the test period? \\
        \textbf{Q15} & Do you have any other comments related to your engagement in Hover? \\ \hline
    \end{tabularx}
    \caption{Questions related to players' engagement.}
    \label{tab:player_engagement_questions}
\end{table}

\subsubsection{Player Motivation}
To investigate the effect game metrics and social interaction features have on the player’s motivation, the participants were asked to respond to nine statements and four questions. The statements were focused on how the different reward systems affected their motivation to do activities. Table \ref{tab:player_motivation_statements} present the statements related to player motivation.

\begin{table}[H]
    \centering
    \begin{tabularx}{0.95\textwidth}{|c|X|}
        \hline
        \textbf{ID} & \multicolumn{1}{c|}{\textbf{Statements}} \\ \hline
        \textbf{S7} & I was motivated to carry out activities by the opportunity to win medals. \\
        \textbf{S8} & I was motivated to carry out activities by the opportunity to earn points. \\
        \textbf{S9} & I was motivated to carry out activities because others did it. \\
        \textbf{S10} & Reactions to my activities inspired me to use the application more. \\
        \textbf{S11} & The leaderboard motivated me to get as many points as possible. \\
        \textbf{S12} & The ability to earn double points by using the application with a friend inspired me to use it more. \\
        \textbf{S13} & Challenges between friends inspired me to use the application. \\
        \textbf{S14} & I was inspired to use Hover by other players' activities. \\
        \textbf{S15} & I opened the application when I received a notification. \\ \hline
    \end{tabularx}
    \caption{Statements related to players' motivation.}
    \label{tab:player_motivation_statements}
\end{table}

Furthermore, the participants were asked to quantify how often they felt inspired to do activities from the game’s different features. For questions Q15-Q17 the participants were given four options: \textit{Every day}, \textit{sometimes}, \textit{a few times} and \textit{never}. Q18 was more of a feedback question to ensure that the players felt that the time they invested was rewarded appropriately. Q19 gave the participant the option to air out any other remarks they had on their motivation to play. The questions on player motivation are presented in Table \ref{tab:player_motivation_questions}.

\begin{table}[H]
    \centering
    \begin{tabularx}{0.95\textwidth}{|c|X|}
        \hline
        \textbf{ID} & \multicolumn{1}{c|}{\textbf{Questions}} \\ \hline
        \textbf{Q16} & How often during the test period were you inspired by your friends to use Hover? \\
        \textbf{Q17} & How often during the test period were you inspired by challenges to use Hover? \\
        \textbf{Q18} & How often during the test period were you inspired by medals to use Hover? \\
        \textbf{Q19} & Was the number of points in each category correct, in relation to your effort? \\
        \textbf{Q20} & Do you have anything else you want to say related to your motivation to Hover? \\ \hline
    \end{tabularx}
    \caption{Questions related to players' motivation.}
    \label{tab:player_motivation_questions}
\end{table}

\subsubsection{Player Enjoyment}
The participants were given nine statements to gather data on perceived enjoyment. Since the game is designed using game theory of chapters \ref{chap:design_an_enjoyable_game}-\ref{chap:player_types}, the statement provides data to investigate if the game elements have the desirable effect on the player’s enjoyment. Statements S16 and S17 directly address the player’s feelings towards Hover, whereas S18-S24 deals with the perceived enjoyment towards specific game elements that we have prioritized to incorporate in our game. The statements directed towards specific game elements are based on the gameflow concepts clear goal, challenge, feedback, and progress. Table \ref{tab:player_enjoyment_statements} present the statements related to player enjoyment.

\begin{table}[H]
    \centering
    \begin{tabularx}{0.95\textwidth}{|c|X|}
        \hline
        \textbf{ID} & \multicolumn{1}{c|}{\textbf{Statements}} \\ \hline
        \textbf{S16} & I liked the concept of Hover. \\
        \textbf{S17} & It was sad to stop using Hover after the test period ended. \\
        \textbf{S18} & I felt that the progress in the game was always clear. \\
        \textbf{S19} & There was an appropriate degree of difficulty to achieve medals. \\
        \textbf{S20} & My status in the game increased the more I played. \\
        \textbf{S21} & The challenges I participated in had clear goals. \\
        \textbf{S22} & I felt attached to my unique avatar. \\
        \textbf{S23} & I was happy to get reactions to my activities and medals. \\
        \textbf{S24} & I was curious to see what my friends had been up to. \\
         \hline
    \end{tabularx}
    \caption{Statements related to players' enjoyment.}
    \label{tab:player_enjoyment_statements}
\end{table}

Moreover, to collect additional questions about their perceived enjoyment, we asked the participant to air out any other remarks on their perceived enjoyment. The questions on enjoyment are presented in Table \ref{tab:player_enjoyment_questions}.

\begin{table}[H]
    \centering
    \begin{tabularx}{0.95\textwidth}{|c|X|}
        \hline
        \textbf{ID} & \multicolumn{1}{c|}{\textbf{Questions}} \\ \hline
        \textbf{Q21} & Was there anything other than what is included in the questionnaire that you enjoyed or found amusing (or not) from using Hover during the test period? \\
        \textbf{Q22} & Do you have any other comments related to your enjoyment in Hover? \\ \hline
    \end{tabularx}
    \caption{Questions related to players' enjoyment.}
    \label{tab:player_enjoyment_questions}
\end{table}

\subsubsection{Player Habits}
To investigate if and how players changed their habits, we included some statements and questions relating to habits and perceived usefulness. The participants were asked to state whether they had changed their habits and whether they were inspired to do more activities during the two weeks. These statements (see Table \ref{tab:player_habits_statements}) are used to study if the concept has the positive effect on individuals that we seek.

\begin{table}[H]
    \centering
    \begin{tabularx}{0.95\textwidth}{|c|X|}
        \hline
        \textbf{ID} & \multicolumn{1}{c|}{\textbf{Statements}} \\ \hline
        \textbf{S25} & Hover has changed my habits about how often I use cultural offerings. \\
        \textbf{S26} & Hover made me want to use cultural offerings more often.\\
        \textbf{S27} & Hover has changed my habits about how often I exercise. \\
        \textbf{S28} & Hover made me want to exercise more often. \\ 
        \textbf{S29} & Hover has changed my habits about how often I visit my educational institution. \\
        \textbf{S30} & Hover made me want to visit campus more often. \\
        \textbf{S31} & Hover has changed my habits about how often I am social. \\ 
        \textbf{S32} &  Hover made me want to visit social zones more often.\\ \hline
    \end{tabularx}
    \caption{Statements related to habits in Hover.}
    \label{tab:player_habits_statements}
\end{table}

Moreover, we also asked the participants to state whether they would like to continue using a game like this as it reveals how they perceive the usefulness of such a concept. Due to the ongoing pandemic, there were several concerns that it would affect the results or conditions for testing the game. The potential change in conditions was local and national restrictions, such as closed campuses, gyms, and restaurants, as well as canceled cultural happenings. It was also realistic to think that the participants' habits and willingness to do activities outside their home deviated from their pre-pandemic behavior. For that reason, Q23 was included to collect data that potentially could support this underlying concern. The questions on habits and usefulness are presented in Table \ref{tab:player_enjoyment_questions}.

\begin{table}[H]
    \centering
    \begin{tabularx}{0.95\textwidth}{|c|X|}
        \hline
        \textbf{ID} & \multicolumn{1}{c|}{\textbf{Questions}} \\ \hline
        \textbf{Q23} & Would you like to continue using a game like this?  \\
        \textbf{Q24} & Do you think you would have used Hover more if it had not been for an ongoing pandemic? \\
        \textbf{Q25} & Do you have any other comments related to your habits before, during, or after using Hover? \\ \hline
    \end{tabularx}
    \caption{Questions related to habits in Hover.}
    \label{tab:player_habits_questions}
\end{table}

\subsection{Interview}
To collect qualitative data on the user’s perception, we invited some of the subjects to participate in individual interviews and to participate in group interviews. Our interviews were unstructured, and the purpose was to discover new information and triangulate other data sources. Our interviews were unstructured, and the purpose was to discover new information and triangulate other data sources. The combination of individual and group interviews have several advantages. Group interviews are suitable for generating “more responses, and more varied responses, as one participant’s views are challenged by others or stimulate others to new ideas” \citep{oates_researching_2006}. On the other hand, single interviews are reliable as it lowers the barrier for voicing opinions that deviate from the masses. An essential consideration for group interviews is to make sure that the group participants are of similar status so that they feel comfortable with speaking their minds and sharing their opinions and thoughts with the group. This was taken into account when deciding what subjects to invite to a group interview.

The unstructured interviews, individual and group, were focused on the topics related to our research questions. This means that there was no defined set of questions, and the conversation was naturally steered towards the topics of player engagement, motivation, enjoyment, and habits.

To collect data from these interviews, we chose to take field notes and audio recorded the interviews to transcribe later. This method allowed us to structure the interviews as a natural conversation rather than a formal interview setting.


\subsection{In-game Observations}
The observations of the user testing were not systematic observation, nor was it a typical participant observation. We found that the best way to observe the subjects was to participate during the testing passively. This presence allowed us to observe how the test subject played and interacted. The test subjects were informed beforehand of our passive participation. We chose the role of passive observers to not compromise the results. We refrained from creating challenges and reacting to activities so that the players would not feel pressured to play the game.

\subsection{Ethics}
Due to the need for user testing, we focused, throughout the project, on issues that could occur related to privacy. Therefore, it was important to be transparent with the participants with how we collect and present their data in the thesis.
It was a priority to develop the app to store as little as possible about the use. However, for the purpose of the study and the game concept, some information such as locations, activity, and user data was necessary.

As the project is dependent on user testing we started a dialogue with the Norwegian Centre for Research Data (NSD) to ensure that our research was completed in a secure, rightful, and ethical manner. This early adoption of communication resolved potential GDPR and privacy issues early on in the project’s phase.

Before launching the game on App Store and Google Play, we saw the need for creating a privacy policy (see Appendix \ref{appendix:privacy_policy}) required by said platforms. The privacy policy discloses how, why, and what data we collect, and we made it available on our website and in the settings menu inside the application. 
